\documentclass[12pt,a4paper]{article}
\usepackage{mwe}
\usepackage{imakeidx}
\makeindex
\usepackage{hyperref}
\usepackage{xcolor}

\begin{document}
    \begin{titlepage}
    	\centering
        {\textsc{Universitá degli studi di Verona} \par}
        \par\vspace{1cm}
    	\includegraphics[width=0.15\textwidth]{logoUni.png}
        \par
    	\vspace{1cm}
    	{\Large \textsc{Progetto di Big Data}\par}
    	\vspace{1cm}
        {\Large \textsc{Wikipedia inverted index}\par}
    	\vspace{5cm}
    	{\Large\itshape Fabrizio Bassan VR479096\par}
    	\vfill
    	{\large Aprile 2023 \par}
    \end{titlepage}
    
    \tableofcontents
    \newpage
    \section{Introduzione}
    Richiesta del progetto: "L’output deve contenere, per ogni parola, la lista di pagine di wikipedia che contengono quella parola".
    Divideremo, quindi, in 3 parti il progetto: una fase di processing, una parte per la gestione dell'inverted index e la fase finale, la creazione dei grafici.
    Utilizzeremo principalmente 4 librerie: pyspark, lxml, pandas, matplotlib.
    Ce ne sono altre ma sono tendenzialmente di contorno a queste che sono praticamente essenziali.
    \newline
    Ovviamente ci sará una prima fase per controllare se il file (con estensione xml) esiste, in quel caso si andrá avanti, altrimenti verrá generato un errore che bloccherá l'esecuzione.
    \section{Processing}
    La fase di processing si suddivide in piú parti.
    \subsection{Metodi}
    Qui abbiamo utilizzato due metodi (removeNestedParentheses e cleanUp) per pulire il piú possibile il file che verrá generato a partire dal file.xml in quanto questo é scritto in formato \textit{WikiText}, un linguaggio di MarkUp particolare usato da wikipedia. Arriviamo ora alle stop words. Cosa sono le stop words? Chiedendo aiuto ad internet, la sua definizione risulta questa: \textit{sono quelle parole che, per la loro alta frequenza in una lingua, sono di solito ritenute poco significative dai motori di ricerca, che le ignorano.}. Nel nostro caso, prendiamo le stop words inglesi, come ad esempio verbi ausiliari (is, are, am, ecc.), congiunzioni (but, and, ecc.) e cosí via.
    \subsection{Analisi del file}
    Finalmente arriviamo alla parte piú importante: il processamento del file.
    \newline
    Creo due file CSV (cosí mi é piú semplice analizzarli piú avanti), con due header ciascuno: index, che é la colonna comune fra i due, e altri, complessivamente, due campi ossia \textbf{text} e \textbf{title}. Molto importante sará l'utilizzo del comando \textit{\textcolor{blue}{elem.clear()}} che ci permetterá di non sovraccaricare la memoria primaria scrivendo quindi ogni volta che un tag \textcolor{red}{text} oppure \textcolor{red}{title} viene trovato, una riga in uno dei due file csv creati all'inizio della cella.
    \newline
    Ho anche impostato un contatore formattato appositamente che stamperá la percentuale di completamento dell'operazione.
    \section{Inverted Index}
    Qui si andrá ad utilizzare la libreria PySpark. Dopo aver creato una SparkSession, andró a unire i due file csv precedentemente creati utilizzando la colonna in comune: \textit{index}. Ne uscirá quindi un file unico contenente, per ogni titolo, tutte le parole che aveva all'interno della sua pagina. Successivamente si andrá a creare un altro file (countedWords) che conterrá 3 colonne: il titolo della pagina, la parola singola utilizzata ed un counter per tener conto di quante volte quella parola é stata utilizzata in quella pagina.
    Da qui, ho generato un altro file (singleWordCounted) che avrá due sole colonne che indicheranno la parola ed il numero di occorrenze.
    Riprendendo \textcolor{blue}{countedWords.csv} vado ad eliminare la colonna \textit{count} cosí da creare un file (finalFileToAnalyze) per lettere tutte le parole presenti e salvarle in un set (per eliminare tutti i doppioni) e fare un casting di questo set a lista in modo da renderlo piú flessibile e comodo.
    \newline
    Ci servirá per creare (finalmente) il nostro file finale: results.csv che avrá due colonne: word e titles (che é la lista dei titoli), filtrando le parole in modo da evitare doppioni, usando la lista precedentemente creata e scrivendo, sempre in un file csv, riga per riga.
    \newline
    Poi creeró altri tre files, initialLetter(due colonne, char e occurence), titlesSize (due colonne, word e list\_size) e occurences (list\_size e occurences che contiene la frequenza del numero di parole) che ci serviranno per la generazione dei grafici.
    \section{Grafici}
    Le librerie che ho utilizzato di piú qui sono principalmente pandas e matplotlib.
    Ho creato 5 grafici:
    
\end{document}